\section{The residue-twisted Cayley character}

\subsection{The quotient and the projective action}

Set
\[
 \mathcal F=yC+zQ,\qquad
 R=K[x_0,\ldots,x_7,y,z]/J(\mathcal F),
\]
with bidegrees
\[
 \deg x_i=(0,1),\qquad
 \deg y=(1,-3),\qquad
 \deg z=(1,-2).                                    \tag{5.1}
\]
The Cayley description \cite[\S5.4, equation~(8)]{FIK2012} gives
\[
 H^{4,1}_{\mathrm{prim}}(X)\cong R_{1,-3},
 \qquad
 H^{3,2}_{\mathrm{prim}}(X)\cong R_{2,-3}.         \tag{5.2}
\]

Choose a lift \(M_g\in\operatorname{GL}_8(K)\) and let \(A_g\) be
defined by
\[
 \binom C Q(M_gx)=A_g\binom C Q(x).                \tag{5.3}
\]
The induced Cayley transformation is
\[
 (x,(y,z)^t)\longmapsto
 (M_gx,A_g^{-t}(y,z)^t).                           \tag{5.4}
\]
The residue class has the additional orientation multiplier
\begin{equation}
\label{eq:orientation}
 \omega(g)=\frac{\det M_g}{\det A_g}.
\end{equation}

\begin{lemma}
\label{lem:lift}
The action consisting of polynomial pullback under \((5.4)\) and the
multiplier \(\omega(g)\) in \eqref{eq:orientation} is independent of
the scalar lift \(M_g\).
\end{lemma}

\begin{proof}
Replacing \(M_g\) by \(\lambda M_g\) replaces \(A_g\) by
\(\operatorname{diag}(\lambda^3,\lambda^2)A_g\).
For a representative in \(R_{p,-3}\), the simultaneous Cayley
polynomial pullback has total scalar \(\lambda^{-3}\): the \(x\)-,
\(y\)-, and \(z\)-degrees combine using the second bidegree \(-3\).
The determinant ratio changes by
\(\lambda^{8-3-2}=\lambda^3\), so the factors cancel.  For the explicit
control \(\lambda=2\), these are \(2^{-3}\) and \(2^3\).
\end{proof}

For the chosen generators,
\[
 A_r=A_s=\operatorname{diag}(1,\rho),\qquad
 \det M_r=\det M_s=\rho,
\]
so \(\omega(r)=\omega(s)=1\).  This accidental simplification does not
justify removing \eqref{eq:orientation}.

\subsection{Exact character}

The ambient bidegree-\((2,-3)\) monomials form
\[
 y^2K[x]_3\oplus yzK[x]_2\oplus z^2K[x]_1,
\]
of dimension
\[
 \binom{10}{3}+\binom92+8=164.
\]
The corresponding Jacobian relations have rank \(81\), hence
\(\dim R_{2,-3}=83\).  Exact quotient matrices give
\begin{align}
 \tr(r^k)
 &=(83,-1,-3,-1,-7,-1,3,-1,-7,-1,-3,-1),
                                                        \tag{5.5}\\
 \tr(sr^k)&=3\qquad(0\le k\le11).                       \tag{5.6}
\end{align}

The four one-dimensional characters, ordered by
\[
 (r,s)\longmapsto(1,1),(1,-1),(-1,1),(-1,-1),
\]
occur with multiplicities
\[
 (4,1,3,3).                                        \tag{5.7}
\]
For \(1\le j\le5\), let
\[
 \vartheta_j(r^k)=2\cos\left(\frac{\pi jk}{6}\right),
 \qquad \vartheta_j(sr^k)=0.
\]
Their multiplicities are
\[
 (7,8,6,8,7).                                      \tag{5.8}
\]
The dimension check is
\[
 4+1+3+3+2(7+8+6+8+7)=83.
\]

Finally, \(R_{1,-3}=Ky\) is the trivial representation.  The group
average therefore retains one dimension in \(H^{4,1}\), four in
\(H^{3,2}\), and their conjugates.  This proves the first Hodge ledger
in Theorem~\ref{thm:main}.  Subtracting from the total
\((1,83,83,1)\) gives the second.

\begin{remark}
The release computation must reconstruct the quotient action and the
character over an exact coefficient field.  A trace list supplied by
the producer without an independent quotient replay is not sufficient
evidence.
\end{remark}
